% Préambule du livre. Assemblé par tools/generate-book.py avec les documents de
% chapitre et les textes de liaison de book/textes/.
%
% La mise en page vise celle d'un livre technique : format 17 × 24 cm, texte en
% serif et titres en sans-serif, encadrés à icônes pour les notes et les mises
% en garde. Les documents de chapitre, eux, gardent leur composition d'article :
% ils se lisent seuls, à côté de leur fichier Lean, et n'ont pas à faire livre.
%
% Le document se compile avec XeTeX (donc avec tectonic), et non plus avec
% pdfLaTeX : fontspec et awesomebox l'exigent.
\documentclass[11pt,DIV=12,BCOR=6mm,open=any,numbers=noenddot,chapterprefix=false]{scrbook}

% Format 17 × 24 cm, celui des livres scientifiques français.
\usepackage[paperwidth=17cm,paperheight=24cm,
  inner=2.1cm,outer=1.9cm,top=2.2cm,bottom=2.4cm]{geometry}

% Le texte en XCharter — un Charter, lisible en petit corps — et les titres en
% Fira Sans. L'ordre compte : les paquets de polices avant fontspec.
\usepackage{XCharter}
\usepackage[scaled=0.94]{FiraSans}
\usepackage{fontspec}
% La chasse fixe porte les identifiants Lean, pleins d'exposants et d'indices
% Unicode (`aᵐ⁺ⁿ`, `uₙ`) : Latin Modern Mono ne les a pas et les laisse tomber
% sans bruit. DejaVu les couvre tous, et sa chasse réduite fait tenir les noms
% longs dans une page de 17 cm.
\setmonofont{DejaVuSansMono.ttf}[Scale=0.82]
\usepackage[french]{babel}
\usepackage{amsmath,amssymb,amsthm}
\usepackage{microtype}

% Les encadrés à icônes : note, astuce, mise en garde, avertissement.
\usepackage{awesomebox}
\usepackage[most]{tcolorbox}

% Les figures : TikZ, pour que leurs étiquettes soient composées dans les
% polices du document. Voir la skill `illustrate-theorem`.
\usepackage{tikz}
\usetikzlibrary{calc}

% Un filigrane sur chaque page : ce livre est une expérimentation de
% formalisation, pas un manuel. Le dire une fois dans l'avant-propos ne suffit
% pas — une page imprimée circule seule.
\usepackage{draftwatermark}
\SetWatermarkText{Expérimentation \textemdash{} ne pas utiliser en classe}
\SetWatermarkScale{0.24}
\SetWatermarkColor[gray]{0.9}
\SetWatermarkAngle{45}

\usepackage[hidelinks]{hyperref}

% Les identifiants Lean en machine à écrire ne se coupent pas : mieux vaut des
% espacements plus lâches que des lignes qui débordent dans la marge.
\emergencystretch=2em

% La table des matières : un numéro comme « 11.2.10 » déborde de la largeur que
% KOMA réserve par défaut, et vient coller le titre. On élargit la colonne des
% numéros à mesure qu'on descend dans la hiérarchie.
\DeclareTOCStyleEntry[numwidth=2.4em]{tocline}{section}
\DeclareTOCStyleEntry[numwidth=3.6em]{tocline}{subsection}
\DeclareTOCStyleEntry[numwidth=4.4em]{tocline}{subsubsection}

% Titres et en-têtes en sans-serif, comme dans un manuel technique.
\setkomafont{disposition}{\sffamily}
\setkomafont{descriptionlabel}{\sffamily\bfseries}
\addtokomafont{caption}{\small}
\usepackage[automark]{scrlayer-scrpage}
\pagestyle{scrheadings}
\clearpairofpagestyles
\automark[section]{chapter}
\ihead*{\small\sffamily\headmark}
\ohead*{\small\sffamily\pagemark}

% Numérotation par chapitre : « Théorème 4.3 » se retrouve dans un livre de
% trois cents pages, « Théorème 137 » non.
\theoremstyle{definition}
\newtheorem{definition}{Définition}[chapter]
\newtheorem{exemple}[definition]{Exemple}
\theoremstyle{plain}
\newtheorem{theoreme}[definition]{Théorème}
\newtheorem{lemme}[definition]{Lemme}

% Un exemple ne se démontre pas : on explique ce qu'il montre.
\newenvironment{explication}{\begin{proof}[Explication]}{\end{proof}}

% Un exemple se donne à voir : un cas calculé, un contre-exemple. Il est composé
% en bloc détaché, filet et fond tramés, comme les encadrés d'un livre technique
% \textemdash{} et son explication vient dans le même bloc, puisqu'elle ne vaut
% que pour lui.
\definecolor{tonexemple}{HTML}{B45309}
\tcolorboxenvironment{exemple}{
  enhanced, breakable, sharp corners, boxrule=0pt,
  colback=tonexemple!4, borderline west={2.5pt}{0pt}{tonexemple!70},
  left=10pt, right=8pt, top=6pt, bottom=6pt,
  before skip=10pt, after skip=2pt,
}
\tcolorboxenvironment{explication}{
  enhanced, breakable, sharp corners, boxrule=0pt,
  colback=tonexemple!4, borderline west={2.5pt}{0pt}{tonexemple!70},
  left=10pt, right=8pt, top=0pt, bottom=6pt,
  before skip=0pt, after skip=10pt,
}

% Un énoncé écrit mais non démontré : la preuve Lean est un `sorry`. La mise en
% garde à icône dit ce qu'un paragraphe en italique disait plus discrètement.
\newcommand{\admis}{\warningbox{\textbf{Admis.} L'énoncé est écrit en Lean, sa
démonstration reste à faire.}}

% Renvoi à la déclaration Lean, une ligne après la démonstration, aligné à
% droite. C'est le seul lien entre le livre et le dépôt : il permet d'aller
% vérifier, ligne par ligne, que la démonstration lue est bien celle que la
% machine accepte.
\newcommand{\source}[2]{\par\smallskip\noindent\hfill{\footnotesize\href{#1}{\texttt{#2}}}\par\smallskip}

% Les repères de l'index. Ils sont posés par generate-book.py sur les termes
% définis et les titres de section, et relus par \pageref à la fin du livre :
% aucun outil extérieur n'est nécessaire, deux passes de compilation suffisent.
\newcommand{\reperage}[1]{\phantomsection\label{#1}}

\title{%
  Mathématiques du secondaire\\
  \large démontrées en Lean 4, de la sixième à la terminale}
\author{Michel Hua \\ \texttt{commutator.io} \\[8pt]
  \normalsize\normalfont assisté d'Opus~5 d'Anthropic\\
  \normalsize\normalfont et des outils de la \emph{Lean Community}}
\date{}
